% P7 student report starter. Replace every [TODO: ...] before submission.
% Coordinate authors, content, citation and timing with the organizers first.
% Never add challenge data, answer CSVs, gold, private evaluation details or
% organizer files to the manuscript, source archive or ancillary attachments.
\documentclass[11pt,a4paper]{article}
\usepackage[T1]{fontenc}
\usepackage[utf8]{inputenc}
\usepackage{lmodern}
\usepackage[margin=25mm]{geometry}
\usepackage{amsmath,booktabs,tabularx,enumitem}
\usepackage{xurl}
\usepackage[hidelinks]{hyperref}
\urlstyle{same}
\setlength{\parindent}{0pt}
\setlength{\parskip}{6pt}
\setlength{\emergencystretch}{2em}
\setlist{nosep,leftmargin=*}
\newcommand{\todo}[1]{\textit{[TODO: #1]}}

\title{\todo{Your method title}\\
  \large SIXT + TUM Hackathon 2026 --- Problem 7}
\author{\todo{Student author names}\\
  \small \todo{Actual affiliations and contact email}}
\date{\todo{Manuscript date}}

\begin{document}
\maketitle

\begin{abstract}
\todo{In about 150 words, state the task, your method, which work models and
tools perform, measured results, and the main limitation. Do not invent
measurements or claim official scores before receiving them.}
\end{abstract}

\section{Task and data}
We study Problem 7 of the SIXT + TUM Hackathon: calculate answers from
synthetic enterprise tables and business rules, then submit one scalar answer
per question in a two-column CSV. The challenge uses the Enterprise Data Agent
Challenge 2026 dataset~\cite{tzachristas2026enterprise}. There are two separate
snapshots, with 25 CSV data tables and 200 scored questions per stage, plus
10 unscored practice examples in the initial stage.

\todo{Explain the business problem and the main calculation challenges in your
own words. Describe your use of the supplied rules and metadata. Keep the two
snapshots separate. Illustrate with public toy examples only.}

\subsection{Dataset identification}
The bibliography contains the prepared Version 6.0.0 citation. Confirm it
against the citation supplied with the issued package. Record each issued
archive below; dataset version and event content revision are separate fields.
Use the organizers' hash for the exact archive you received.

% Repeat complete fields for both stages. Use \nolinkurl{...} for filenames,
% hashes and revisions so underscores and long identifiers wrap safely.
\begin{description}[style=nextline]
  \item[Initial/public stage]
    Archive: \todo{filename}; SHA-256: \todo{64-character hash};
    access date: \todo{YYYY-MM-DD}; dataset version: \todo{issued version};
    event content version: \todo{issued revision};
    event policy version: \todo{issued policy}.
  \item[Final stage]
    Archive: \todo{filename}; SHA-256: \todo{64-character hash};
    access date: \todo{YYYY-MM-DD}; dataset version: \todo{issued version};
    event content version: \todo{issued revision};
    event policy version: \todo{issued policy}.
\end{description}

\section{Method}
\todo{Describe the workflow from a question to the exported answer. Explain
table selection, joins, record eligibility, dates, currencies, rounding,
validation and human intervention. State what models do and what Python,
SQL or other deterministic tools do. Justify the design and discuss a baseline
where available. A simple diagram of your own workflow is welcome.}

\subsection{Implementation and reproducibility}
\todo{Give the framework/runtime, dependencies, setup and execution commands,
model settings, random seeds where relevant, and code revision used for each
stage. Describe the execution environment. Provide a code link only if it is
authorized for public sharing and contains no restricted material or secrets.
Explain how organizers with the issued data can reproduce the workflow.}

\section{Evaluation and results}
\todo{Describe your validation procedure and distinguish your own checks from
organizer-reported results. Report only measurements you actually obtained.
Use ``not measured'' or ``not available'' for missing values. Discuss errors
without reproducing competition questions or private answer-level feedback.}

\begin{table}[htbp]
\centering
\caption{Measured results. Replace placeholders; give score provenance and
measurement conditions in the text.}
\begin{tabularx}{\linewidth}{@{}Xll@{}}
\toprule
Measure & Public stage & Final stage \\
\midrule
Official correct answers (out of 200) & \todo{value} & \todo{value} \\
Official accuracy (\%) & \todo{value} & \todo{value} \\
Elapsed workflow time & \todo{value} & \todo{value} \\
Model tokens used & \todo{value} & \todo{value} \\
Measured model/API cost (currency) & \todo{value} & \todo{value} \\
\bottomrule
\end{tabularx}
\end{table}

Report hardware, memory use and the scope of timing/cost measurements if
measured. Do not confuse accuracy with the Hackathon grade. For reference,
with public and final accuracies $P,F$ expressed from 0 to 100,
\[
 D=0.35P+0.65F,\qquad G=0.65D+J+B,
\]
where $J$ is the judges' contribution (0--35 points) and $B$ is the verified
small-model bonus (0 or 15 points). Thus $G$ is at most 115 points.
\todo{If reporting a grade, use the organizers' confirmed components.}

\section{Model-use declaration}
\textbf{Complete this section for every team, whether or not claiming the
bonus.} Reconcile it with the submitted \texttt{MODEL\_DECLARATION.csv}.
Disclose all generative models used for challenge-specific design, coding,
planning, analysis, answer generation or validation across both stages,
including coding assistants, validators and fallback models. Also disclose
any additional model assistance used to prepare the later manuscript,
clearly separated from the challenge workflow.

% Duplicate this block once per model. Never put credentials in the report.
\subsection*{Model \todo{number}: \todo{exact identifier}}
\begin{description}[style=nextline]
  \item[Identity and size] Publisher: \todo{publisher}; exact version/revision:
    \todo{version}; total parameters in billions: \todo{count or unknown};
    parameter-count source: \todo{publisher/model-card reference and URL}.
  \item[Role and stage] \todo{Specific task, initial/final/both stages,
    substantive contribution, coding or fallback use; distinguish any
    manuscript-only use after the Hackathon.}
  \item[Configuration] \todo{Framework/runtime and version, local or hosted,
    quantization, decoding settings, prompts/configuration version and tools.}
  \item[Usage and evidence] \todo{Measured usage or not measured;
    evidence supplied to organizers; a sanitized public summary where allowed.
    Do not embed private prompts, data rows, answers, logs or secrets.}
\end{description}

\subsection*{Small-model bonus declaration}
\begin{itemize}
  \item Bonus claimed for the Hackathon workflow: \todo{yes/no}.
  \item Every generative model in the challenge workflow has at most
    40 billion total parameters: \todo{yes/no/unverifiable}.
  \item Meaningful analytical work performed by qualifying models:
    \todo{explain, or not applicable}.
  \item Size and usage evidence supplied to organizers:
    \todo{identify evidence, without publishing restricted material}.
  \item Organizer verification outcome: \todo{pending/confirmed/ineligible}.
\end{itemize}
For mixture-of-experts models, count all parameters, including inactive
experts. Quantization does not reduce the parameter count. Each generative
model must meet the limit; counts are not added across models. Unknown sizes
cannot establish eligibility. A model-free workflow does not earn the bonus.
The declaration alone does not establish eligibility; the published P7
policy and organizer verification apply.

\section{Limitations and business relevance}
\todo{Discuss failure cases, reliability, human effort, reproducibility,
efficiency trade-offs and what would change with more time. Explain how the
approach could be reused in a SIXT-like business setting. Distinguish evidence
from speculation. Do not imply validation on real SIXT data.}

\section*{Author contributions and acknowledgements}
\todo{State each student's contribution, acknowledge assistance, cite reused
code, models and frameworks, and give accurate funding or conflict disclosures
if applicable. Confirm the author list with coauthors and organizers.}

\section*{Data and code availability}
The challenge data were supplied separately for the Hackathon. This report
does not redistribute the datasets, answer CSVs, golden answers, private
evaluation details or organizer material.
\todo{State the actual availability of your method code. Include only
authorized public links. Do not imply an open dataset license or invent a DOI.
If organizers later provide a verified dataset repository citation, update
the bibliography accordingly.}

\bibliographystyle{plain}
\bibliography{references}
\end{document}
